\documentclass{article}
\usepackage[left=2cm, right=2cm, top=2cm, bottom=2cm]{geometry}
\usepackage{graphicx}
\usepackage{amsmath}
\usepackage{amssymb}
\usepackage{amsthm}
\usepackage{fancyhdr}
\usepackage{verbatim}
\usepackage{listings}
\usepackage{xcolor}
\usepackage{pdfpages}
\usepackage{cancel}
\usepackage{physics}
\usepackage{esint}
\usepackage{braket}

\lstset{
    language=C++,
    basicstyle=\ttfamily\footnotesize,
    keywordstyle=\color{blue},
    commentstyle=\color{green},
    stringstyle=\color{red},
    numbers=left,
    numberstyle=\tiny\color{gray},
    frame=single,
    breaklines=true,
    captionpos=b,
}


\title{PHYS 427: Lecture \# 2}
\author{Cliff Sun}

\newtheorem{theorem}{Theorem}[section]
\newtheorem{lemma}[theorem]{Lemma}
\newtheorem{definition}[theorem]{Definition}
\newtheorem{conjecture}[theorem]{Conjecture}
\newtheorem{proposition}[theorem]{Proposition}
\newtheorem{corollary}[theorem]{Corollary}
\newtheorem{one minute paper}[theorem]{One Minute Paper}

\pagestyle{fancy}
\lhead{\textbf{\thepage}\ \ \nouppercase{\rightmark}}
\chead{PHYS 427: Lecture \# 2}
\rhead{Cliff Sun}

\begin{document}

\maketitle

\section*{Lecture Span}
\begin{enumerate}
    \item Multiplicity of Physical Systems
\end{enumerate}

\section*{Review}

Useful: 

\begin{equation}
    (p + q)^N = \sum^N_{n=0}\begin{pmatrix}
        N \\
        n
    \end{pmatrix}p^nq^{N-n}
\end{equation}

One can also take the following derivative:

\begin{equation}
    p\partial_p p^{n} = n p^n
\end{equation}

Then, taking the derivative of both sides of Eq.(1) and multiply it by $p$, then we obtain:

\begin{align}
    Np(p + q)^{N-1} &= \sum^{N}\begin{pmatrix}
        N \\
        n
    \end{pmatrix}np^n q^{N -n}\\
    Np &= \sum^{N}\begin{pmatrix}
        N \\
        n
    \end{pmatrix}np^n q^{N -n}
\end{align}

Therefore, $\langle n \rangle = pN$. Such a procedure can be repeated for $\langle n^2 \rangle$. Here, 

\begin{align}
    \langle n^2 \rangle &= \sum n^2 \begin{pmatrix}
        N \\
        n
    \end{pmatrix}p^n q^{N-n}
\end{align}

Moreover, 

\begin{equation}
    \left( p \partial_p \right)^2 p^{n} = n^2 p^n
\end{equation}

Then, 

\begin{align}
    \sum \begin{pmatrix}
        N \\
        n
    \end{pmatrix}n^2 p^n q^{N-n} &= \left( p \partial_p \right)^2 (p + q)^N \\
    &= p \partial_p Np(p + q)^{N-1} \\
    &= Np(p + q)^{N-1} + N(N-1)p^2(p +q)^{N-2} \\
    \langle n^2 \rangle &= Np + N(N-1)p^2
\end{align}

Here, $P(n)$ is a peak. The width of that peak is shown as 

\begin{align}
    \sigma_n^2 &= \langle (n - \langle n \rangle)^2 \rangle \\
    &= \langle n^2 \rangle  - \langle n \rangle^2
\end{align}

The relative width can be calculated as $\sigma_n / \langle n \rangle$, and is useful for comparing the true standard deviation versus the magnitude of the atoms we are dealing with.
E.g., a $\sigma_n \sim 1$ and $\mu \sim 10^{23}$ means that $\sigma_n$ doesn't really mean anything. This is essentially a delta function. 

\section*{Multiplicity of a physical system}

We assume that the system is closed (i.e., an isolated system). This also means that the system does not interact with the outside environment. Then, let $U$ be the total energy of the system. Then,

\begin{center}
    U is constant in the isolated system
\end{center}

Later, $U$ can exchange with some sort of bath (environment). The macrostate of the system is defined as $\Omega(U,N)$ where $N$ is the total number of atoms. 

\subsection*{Einstein model of solid}

Assume that we have a 2d lattice of atoms, and each atom are coupled to each other by springs. Then, treat all the atoms as a collection of $N$, independent quantum harmonic oscillators. As a reminder, the energy levels are 

\begin{equation}
    E_i = \hbar \omega (n_i + 1/2), \quad n_i = 0,1,2,\dots
\end{equation}

Here, $i$ is the index of each atom. Then, 

\begin{equation}
    U = \sum_i E_i = \hbar \omega (n + \frac{N}{2})
\end{equation}

Here, we neglect the zero point energy since $\hbar \omega N/2$. 
Moreover, we also have that: $$n = \sum_i^N s_i$$

So then, what is $\Omega(U, N)$? In other words, how many microstates correspond to the same total energy $U$ correspond to $n$ quanta? This is possible when

\begin{equation}
    s_1 + s_2 + \dots = n
\end{equation}

This is also the integer partitions of $n$. We can imagine drawing out $n$ different dots (up to $n$ quanta) and dividing them up. There is a maximum of $N-1$ partitions. 

Then, there are a total of $n + (N-1)$ slots. Then, the bionomial coefficient is 

\begin{equation}
    \begin{pmatrix}
        n + (N-1) \\
        n
    \end{pmatrix}
\end{equation}

Here, $n$ is the number of quanta and $N-1$ is the number of partitions between the atoms. Therefore, there is a total number of $n + N-1$ arrangements of the partitions and the quanta. We also choose $n$ quanta, which is to say how many times 
are we able to arrange $n$ in these slots given that the partitions remain "still". 

We can also approximate 

\begin{equation}
    \Omega(n,N) = \frac{(n + N -1)!}{n!(N-1)!}
\end{equation}

Using stirling's approximation for $N >>1$, we derive that 

\begin{align}
    \ln(N!) &= \ln\left[ N(N-1)(N-2)\dots \right] \\
    &= \ln N + \ln (N-1) + \dots\\
    &= \sum_x \ln x \approx \int_1^N \ln(x)dx\\
    &\approx N \ln N - N
\end{align}

In an Einstein solid, we have that 

\begin{align}
    \ln \Omega(n,N) &= \ln\left[ (n + N)! \right] - \ln \left[N! \right] = \ln \left[ n! \right] \\
    &\approx (n + N)\ln(n + N) - N\ln N - n \ln n
\end{align}

In the high temperature limit, for $n \gg N$, then 

\begin{align}
    \ln (n + N) &= \ln \left( n \left( 1 + \frac{N}{n} \right) \right) \\\
    &= \ln n + \ln \left( 1 + \frac{N}{n} \right) \\
    &\approx \ln n + \frac{N}{n}
\end{align}

Therefore, 

\begin{equation}
    \ln \Omega (n,N) \approx N\ln \left( \frac{n}{N} \right) + N
\end{equation}

Therefore, 

\begin{align}
    \Omega &= e^{\ln \Omega} \approx e^{N \ln (N/n) + N}\\
    &= \left( e \frac{n}{N} \right)^N\\
    &= \left( \frac{eU}{N \hbar \omega} \right)^N
\end{align}

Here, $n = U/\hbar \omega$. Therefore, 

\begin{equation}
    \Omega \sim U^N
\end{equation}

For a lot of systems, we have that 

\begin{equation}
    \Omega \sim U^f
\end{equation}

In an Einstein solid, $f \sim N$. 

\subsection*{Entropy}

Here, we define Entropy as:

\begin{equation}
    \sigma = \ln \Omega
\end{equation}

This is so that $\sigma$ is additive when you add two systems together and calculate its multiplied multiplicity. Then, 

\begin{equation}
    \sigma(U,N) = N \ln \left( \frac{eU}{N \hbar \omega} \right)
\end{equation}

If you double the system, then $U \rightarrow 2U$ and $N \rightarrow 2N$, then $\sigma \rightarrow 2\sigma$. Note, this is only true in the high temperature limit. 

\subsection*{Paramagnetic}

Collection of non-interacting spins. Then, $$N = n_{\uparrow} + n_{\downarrow}$$

Such spins have a dipole moment

\begin{equation}
    \vec{mu} = g \frac{q}{2m}\vec{s}
\end{equation}

Then, the energy is then: (for a B field in the z direction)

\begin{equation}
    U = -\vec{\mu} \cdot \vec{B} = -\mu_z B
\end{equation}

Such that 

\begin{equation}
    \mu_z = g \frac{q\hbar}{2m}m_z, \quad m_z = -s, -s + 1, \dots, s-1, s
\end{equation}

Here, $s = 1/2$ (i.e., we only deal with particles with a spin-1/2.) There are only two energy levels. 

Working within the binary system for $s = 1/2$, then 

\begin{equation}
    \Omega(n_{\uparrow}, N) = \frac{N!}{(n_{\uparrow}!)(N - n_{\uparrow})!}
\end{equation}

Then 

\begin{equation}
    \Omega(U, N) = \frac{N!}{(N/2 - U/(2 \mu B))!(N/2 + U/(2\mu B))!}
\end{equation}

Here, we used the fact that $U = n_{\uparrow}(-\mu B) + n_{\downarrow}(\mu B)$. Here, 

\begin{equation}
    n_{\uparrow} = \frac{N}{2} - \frac{U}{2\mu B}
\end{equation}

\textbf{The above equation is necessary for the homework.} We can also then use Stirling's approximation to then obtain 

\begin{equation}
    \ln \Omega(n_{\uparrow}, N) \iff \sigma (n_{\uparrow, N}) = N \ln N - (N - n_{\uparrow})\ln (N - n_{\uparrow}) - n_{\uparrow}\ln n_{\uparrow}
\end{equation}

You can also re-express this using equation (40) to get $\sigma(U,N)$. $\sigma(U,N)$ as a function of $U$ looks like a half parabola with a maximum at $U = 0$. 

\end{document}